\documentclass{beamer}

% Canonical Beamer model
\usepackage{bookmark}
\usetheme{Madrid}
\usecolortheme{default}
\usepackage{listings}
\usepackage{xcolor}

\lstdefinestyle{code}{
  language=Python,
  basicstyle=\ttfamily\small,
  keywordstyle=\color{blue},
  commentstyle=\color{gray},
  stringstyle=\color{red!60!black},
  showstringspaces=false,
  tabsize=2,
  breaklines=true
}

\setbeamercolor{palette primary}{bg=red, fg=white}
\setbeamercolor{palette secondary}{bg=red!95!black, fg=white}
\setbeamercolor{palette tertiary}{bg=red!90!black, fg=white}
\setbeamercolor{frametitle}{bg=red, fg=white}
\setbeamercolor{title}{bg=red, fg=white}
\setbeamercolor{section in toc}{fg=red}

% Standard packages
\usepackage[utf8]{inputenc}
\usepackage{amsmath}
\usepackage{amssymb}
\usepackage{graphicx}
\usepackage{booktabs}
\usepackage{hyperref}
\usepackage{tikz}
\usetikzlibrary{positioning,arrows.meta,shapes}

% R-specific listing style. Keep R syntax while inheriting the canonical
% visual language for code blocks.
\lstdefinestyle{Rstyle}{
  style=code,
  language=R,
  numbers=left,
  numberstyle=\tiny\color{gray},
  frame=single,
  rulecolor=\color{gray!30},
  columns=fullflexible,
  keepspaces=true
}
\lstset{style=Rstyle}

% Title information
\title[Programming with R]{Programming Fundamentals with R}
\subtitle{Data Structures, Control Flow, Functions, and Data Analysis}
\author{Diogo Ribeiro}
\institute{Faculty of Media Arts and Design, Technical University of Porto}
\date{\today}

\begin{document}

\begin{frame}
  \titlepage
\end{frame}

\begin{frame}{Learning Objectives}
  By the end of this presentation, you should be able to:
  \begin{itemize}
    \item understand R's core data structures and vectorized computation;
    \item write control-flow logic and reusable functions;
    \item manipulate tabular data and perform basic exploratory analysis;
    \item interpret R code with an emphasis on reproducibility and clarity.
  \end{itemize}
\end{frame}

\begin{frame}{Outline}
  \tableofcontents
\end{frame}

\section{Introduction to R}

\begin{frame}{What is R?}
  \begin{itemize}
    \item R is a language and environment for statistical computing and graphics.
    \item It was created by Ross Ihaka and Robert Gentleman in the early 1990s.
    \item R is open source and has a large package ecosystem through CRAN.
    \item It is widely used in statistics, data science, econometrics, and research.
  \end{itemize}
\end{frame}

\begin{frame}{Why Learn R?}
  \begin{itemize}
    \item Strong statistical and visualization libraries.
    \item Interactive workflow for exploratory data analysis.
    \item Excellent support for reproducible research.
    \item Mature ecosystem for specialized scientific domains.
    \item Natural fit for vectorized and matrix-oriented computation.
  \end{itemize}
\end{frame}

\begin{frame}[fragile]{Your First R Commands}
\begin{lstlisting}
# Arithmetic
2 + 3
5 * 4
10 / 2
2^8

# Assignment
x <- 10
y <- 5
z <- x + y

# Output
print(z)
z
\end{lstlisting}
\end{frame}

\section{Basic Data Types}

\begin{frame}{Atomic Data Types}
  Common atomic types in R include:
  \begin{itemize}
    \item \textbf{logical}: \texttt{TRUE}, \texttt{FALSE}
    \item \textbf{integer}: \texttt{1L}, \texttt{10L}
    \item \textbf{double/numeric}: \texttt{3.14}, \texttt{2.0}
    \item \textbf{character}: \texttt{"data science"}
    \item \textbf{complex}: \texttt{1 + 2i}
    \item \textbf{raw}: bytes
  \end{itemize}
\end{frame}

\begin{frame}[fragile]{Inspecting Types}
\begin{lstlisting}
x <- 42
y <- "hello"
z <- TRUE

class(x)
typeof(x)
class(y)
typeof(y)
class(z)
typeof(z)
\end{lstlisting}
\end{frame}

\section{Vectors and Vectorization}

\begin{frame}{Vectors}
  \begin{itemize}
    \item A vector is a one-dimensional collection of values of the same atomic type.
    \item The function \texttt{c()} combines values.
    \item R operations are often vectorized, avoiding explicit loops.
  \end{itemize}
\end{frame}

\begin{frame}[fragile]{Creating Vectors}
\begin{lstlisting}
x <- c(1, 2, 3, 4, 5)

y <- 1:10

z <- seq(from = 0, to = 1, by = 0.1)

repeated <- rep(2, times = 5)

length(x)
\end{lstlisting}
\end{frame}

\begin{frame}[fragile]{Vectorized Operations}
\begin{lstlisting}
x <- c(1, 2, 3, 4, 5)
y <- c(10, 20, 30, 40, 50)

x + y
x * 2
x^2
sqrt(x)

mean(x)
sum(x)
sd(x)
\end{lstlisting}
\end{frame}

\begin{frame}[fragile]{Logical Indexing}
\begin{lstlisting}
x <- c(2, 5, 8, 11, 14)

x > 7
x[x > 7]

# Combine conditions
x[x >= 5 & x <= 11]
\end{lstlisting}
\end{frame}

\section{Matrices and Arrays}

\begin{frame}{Matrices}
  \begin{itemize}
    \item A matrix is a two-dimensional homogeneous data structure.
    \item Matrix algebra is central to many statistical and ML methods.
  \end{itemize}
\end{frame}

\begin{frame}[fragile]{Matrix Basics}
\begin{lstlisting}
A <- matrix(1:6, nrow = 2, ncol = 3)
A

A[1, 2]
A[, 1]
A[2, ]

B <- matrix(7:12, nrow = 3, ncol = 2)
A %*% B
\end{lstlisting}
\end{frame}

\section{Lists and Data Frames}

\begin{frame}{Lists}
  \begin{itemize}
    \item Lists can contain heterogeneous objects.
    \item They are used extensively to return complex results from functions.
  \end{itemize}
\end{frame}

\begin{frame}[fragile]{List Example}
\begin{lstlisting}
student <- list(
  name = "Ana",
  age = 22,
  grades = c(15, 17, 18)
)

student$name
student[[2]]
student$grades
\end{lstlisting}
\end{frame}

\begin{frame}{Data Frames}
  \begin{itemize}
    \item A data frame is a tabular structure.
    \item Each column is a vector; columns may have different types.
    \item Most statistical workflows operate on data frames.
  \end{itemize}
\end{frame}

\begin{frame}[fragile]{Creating a Data Frame}
\begin{lstlisting}
df <- data.frame(
  name = c("Ana", "Bruno", "Carla"),
  age = c(22, 25, 21),
  score = c(17.5, 14.0, 18.5)
)

df
str(df)
summary(df)
\end{lstlisting}
\end{frame}

\begin{frame}[fragile]{Selecting Data}
\begin{lstlisting}
df$name

df[1, ]
df[, "score"]

df[df$score >= 17, ]
\end{lstlisting}
\end{frame}

\section{Factors and Missing Values}

\begin{frame}{Factors}
  \begin{itemize}
    \item Factors represent categorical variables with a finite set of levels.
    \item They are important in statistical modeling, especially for contrasts.
  \end{itemize}
\end{frame}

\begin{frame}[fragile]{Factor Example}
\begin{lstlisting}
group <- factor(c("A", "B", "A", "C", "B"))

group
levels(group)
table(group)
\end{lstlisting}
\end{frame}

\begin{frame}{Missing Values}
  \begin{itemize}
    \item R represents missing values using \texttt{NA}.
    \item Missingness must be handled explicitly in many functions.
    \item \texttt{NaN}, \texttt{Inf}, and \texttt{-Inf} are distinct special values.
  \end{itemize}
\end{frame}

\begin{frame}[fragile]{Working with Missing Data}
\begin{lstlisting}
x <- c(1, 2, NA, 4, 5)

is.na(x)

mean(x)
mean(x, na.rm = TRUE)

x[!is.na(x)]
\end{lstlisting}
\end{frame}

\section{Control Flow}

\begin{frame}{Conditional Logic}
  \begin{itemize}
    \item Use \texttt{if}, \texttt{else if}, and \texttt{else} for branching.
    \item For vectors, prefer vectorized functions such as \texttt{ifelse()} where appropriate.
  \end{itemize}
\end{frame}

\begin{frame}[fragile]{If/Else Example}
\begin{lstlisting}
score <- 16

if (score >= 18) {
  grade <- "Excellent"
} else if (score >= 14) {
  grade <- "Good"
} else if (score >= 10) {
  grade <- "Pass"
} else {
  grade <- "Fail"
}

grade
\end{lstlisting}
\end{frame}

\begin{frame}{Loops}
  R supports:
  \begin{itemize}
    \item \texttt{for}
    \item \texttt{while}
    \item \texttt{repeat}
  \end{itemize}
  Vectorized operations and the apply family are often clearer and faster for many tasks.
\end{frame}

\begin{frame}[fragile]{For Loop}
\begin{lstlisting}
values <- c(2, 4, 6, 8)

for (value in values) {
  print(value^2)
}
\end{lstlisting}
\end{frame}

\begin{frame}[fragile]{While Loop}
\begin{lstlisting}
count <- 1

while (count <= 5) {
  print(count)
  count <- count + 1
}
\end{lstlisting}
\end{frame}

\section{Functions}

\begin{frame}{Why Functions?}
  \begin{itemize}
    \item Reuse logic.
    \item Reduce duplication.
    \item Improve readability and testing.
    \item Separate concerns into small, composable units.
  \end{itemize}
\end{frame}

\begin{frame}[fragile]{Defining Functions}
\begin{lstlisting}
mean_center <- function(x) {
  x - mean(x, na.rm = TRUE)
}

values <- c(2, 4, 6, 8)
mean_center(values)
\end{lstlisting}
\end{frame}

\begin{frame}[fragile]{Default Arguments}
\begin{lstlisting}
power <- function(x, exponent = 2) {
  x^exponent
}

power(4)
power(4, 3)
\end{lstlisting}
\end{frame}

\begin{frame}{Scope}
  \begin{itemize}
    \item R uses lexical scoping.
    \item A function searches for variables in the environment in which it was defined.
    \item Understanding environments becomes important for advanced R programming.
  \end{itemize}
\end{frame}

\section{Apply Family}

\begin{frame}{Functional Iteration}
  Common tools include:
  \begin{itemize}
    \item \texttt{apply()} for matrix margins
    \item \texttt{lapply()} for lists, returning a list
    \item \texttt{sapply()} for simplified output
    \item \texttt{vapply()} for type-safe simplified output
    \item \texttt{tapply()} for grouped summaries
  \end{itemize}
\end{frame}

\begin{frame}[fragile]{Apply Examples}
\begin{lstlisting}
A <- matrix(1:12, nrow = 3)

apply(A, 1, sum)
apply(A, 2, mean)

items <- list(1:3, 4:6, 7:9)
lapply(items, mean)
sapply(items, mean)
\end{lstlisting}
\end{frame}

\section{Basic Data Analysis}

\begin{frame}{A Typical Workflow}
  \begin{enumerate}
    \item Import data.
    \item Inspect structure and quality.
    \item Clean and transform.
    \item Summarize and visualize.
    \item Fit models.
    \item Validate assumptions and interpret results.
    \item Communicate and reproduce the analysis.
  \end{enumerate}
\end{frame}

\begin{frame}[fragile]{Reading CSV Data}
\begin{lstlisting}
# Base R
sales <- read.csv("sales.csv")

head(sales)
str(sales)
summary(sales)
\end{lstlisting}
\end{frame}

\begin{frame}[fragile]{Grouped Summaries with Base R}
\begin{lstlisting}
aggregate(
  score ~ group,
  data = df,
  FUN = mean
)
\end{lstlisting}
\end{frame}

\section{Visualization}

\begin{frame}{Base R Graphics}
  \begin{itemize}
    \item \texttt{plot()}
    \item \texttt{hist()}
    \item \texttt{boxplot()}
    \item \texttt{barplot()}
  \end{itemize}
\end{frame}

\begin{frame}[fragile]{Simple Plot}
\begin{lstlisting}
x <- 1:20
y <- 2 * x + rnorm(20)

plot(
  x, y,
  main = "Simple scatter plot",
  xlab = "x",
  ylab = "y"
)
\end{lstlisting}
\end{frame}

\begin{frame}{The Tidyverse}
  The tidyverse provides a coherent set of packages for data science, including:
  \begin{itemize}
    \item \texttt{dplyr}
    \item \texttt{tidyr}
    \item \texttt{ggplot2}
    \item \texttt{readr}
    \item \texttt{purrr}
    \item \texttt{tibble}
  \end{itemize}
\end{frame}

\begin{frame}[fragile]{A dplyr Pipeline}
\begin{lstlisting}
library(dplyr)

result <- df |>
  filter(score >= 10) |>
  group_by(group) |>
  summarise(
    mean_score = mean(score),
    n = n()
  )
\end{lstlisting}
\end{frame}

\begin{frame}[fragile]{A ggplot2 Example}
\begin{lstlisting}
library(ggplot2)

ggplot(df, aes(x = age, y = score)) +
  geom_point() +
  geom_smooth(method = "lm", se = FALSE) +
  labs(
    title = "Score versus age",
    x = "Age",
    y = "Score"
  )
\end{lstlisting}
\end{frame}

\section{Statistical Modeling}

\begin{frame}{Linear Regression}
  A basic linear model is
  \[
    Y_i = \beta_0 + \beta_1 X_i + \varepsilon_i.
  \]
  In R, linear models are specified using formulas.
\end{frame}

\begin{frame}[fragile]{Fitting a Linear Model}
\begin{lstlisting}
model <- lm(score ~ age, data = df)

summary(model)
coef(model)
fitted(model)
residuals(model)
\end{lstlisting}
\end{frame}

\begin{frame}[fragile]{Prediction}
\begin{lstlisting}
new_students <- data.frame(
  age = c(20, 23, 27)
)

predict(
  model,
  newdata = new_students,
  interval = "confidence"
)
\end{lstlisting}
\end{frame}

\section{Reproducibility and Good Practice}

\begin{frame}{Reproducible R Code}
  \begin{itemize}
    \item Use scripts, notebooks, or Quarto/R Markdown documents.
    \item Record package versions and dependencies.
    \item Use deterministic random seeds where appropriate.
    \item Keep raw data immutable.
    \item Separate data preparation, analysis, and reporting.
    \item Use version control.
  \end{itemize}
\end{frame}

\begin{frame}[fragile]{Random Seeds}
\begin{lstlisting}
set.seed(42)

rnorm(5)
runif(5)
sample(1:20, 5)
\end{lstlisting}
\end{frame}

\begin{frame}{Naming and Style}
  \begin{itemize}
    \item Use descriptive names: \texttt{customer\_age}, not \texttt{x1}.
    \item Prefer small functions with one responsibility.
    \item Be consistent with indentation and spacing.
    \item Comment decisions and assumptions, not obvious syntax.
    \item Avoid hidden global state.
  \end{itemize}
\end{frame}

\section{Exercises}

\begin{frame}{Exercise 1: Vector Operations}
  Given
  \[
    x = (3, 5, 8, 12, 15, 21),
  \]
  write R code to:
  \begin{enumerate}
    \item select values greater than 10;
    \item compute their mean;
    \item standardize the complete vector.
  \end{enumerate}
\end{frame}

\begin{frame}{Exercise 2: Functions}
  Write a function \texttt{summary\_stats(x)} that returns a list containing:
  \begin{itemize}
    \item mean;
    \item median;
    \item standard deviation;
    \item minimum;
    \item maximum.
  \end{itemize}
  Handle missing values correctly.
\end{frame}

\begin{frame}{Exercise 3: Data Frames}
  Create a data frame with at least five observations and columns for:
  \begin{itemize}
    \item name;
    \item group;
    \item age;
    \item score.
  \end{itemize}
  Then compute mean score by group.
\end{frame}

\section{Summary}

\begin{frame}{Key Takeaways}
  \begin{itemize}
    \item R is built around vectors and vectorized operations.
    \item Data frames and lists are fundamental structures for data analysis.
    \item Control flow and functions enable reusable programs.
    \item The apply family and pipelines support concise iteration.
    \item R combines programming, statistics, and visualization in one environment.
    \item Reproducibility and clear code are part of the analytical method.
  \end{itemize}
\end{frame}

\begin{frame}{Where to Go Next}
  \begin{itemize}
    \item Tidyverse data manipulation.
    \item Functional programming with \texttt{purrr}.
    \item Statistical modeling and diagnostics.
    \item Time series and forecasting.
    \item Bayesian analysis.
    \item Package development and testing.
    \item Quarto/R Markdown for reproducible reporting.
  \end{itemize}
\end{frame}

\begin{frame}{Resources}
  \begin{itemize}
    \item R Project: \url{https://www.r-project.org/}
    \item CRAN: \url{https://cran.r-project.org/}
    \item R for Data Science: \url{https://r4ds.hadley.nz/}
    \item Advanced R: \url{https://adv-r.hadley.nz/}
  \end{itemize}
\end{frame}

\begin{frame}
  \centering
  {\Huge Thank You}

  \vspace{0.8cm}
  \textbf{Questions \& Discussion}

  \vspace{1cm}
  \textbf{Diogo Ribeiro}\\
  Faculty of Media Arts and Design, Technical University of Porto\\[0.5cm]
  \texttt{dfr@esmad.ipp.pt}\\
  \texttt{https://orcid.org/0009-0001-2022-7072}
\end{frame}

\end{document}
